\begin{table}[t]\centering\small
\begin{tabular}{llcc}
\toprule
null & $n$ & polynomial block & single index \\
\midrule
additive control & 2000 & 0.02~[0.01--0.05] & 0.10~[0.06--0.14] \\
additive control & 10000 & 0.05~[0.03--0.09] & 0.06~[0.03--0.10] \\
interaction $Z_1Z_2$ & 2000 & \textbf{0.03~[0.01--0.06]} & 0.20~[0.15--0.26] \\
interaction $Z_1Z_2$ & 10000 & 0.05~[0.03--0.09] & 0.09~[0.05--0.13] \\
transcendental $\sin(\textstyle\sum_j Z_j)$ & 2000 & \textbf{1.00~[0.98--1.00]} & 0.15~[0.11--0.21] \\
transcendental $\sin(\textstyle\sum_j Z_j)$ & 10000 & \textbf{1.00~[0.98--1.00]} & 0.06~[0.03--0.10] \\
\bottomrule
\end{tabular}
\caption{Complementary failure of the two nuisance blocks (size, nominal $0.05$; $200$ reps, Wilson 95\% CI; $Z\in\mathbb{R}^3$). The polynomial block removes polynomial interactions at any $n$ but diverges on a transcendental mean; the single index represents the transcendental mean (cleanly once $n$ is large enough to estimate its direction) but underfits the interaction at small $n$. The default test uses both.}
\label{tab:nuisance-tradeoff}
\end{table}
